\section{Satisfiability}
\begin{prob}[Satisfiability Problem]\label{prob:Satisfiability-1}%Problem/Question
Let $\left\lbrace x_i \right\rbrace_{i=1}^{n} $ be boolean variables, and $t_i \in \left\lbrace 
x_i\right\rbrace_{i=1}^{n}  \cup \left\lbrace \bar{x}_i\right\rbrace_{i=1}^{n}   $.  A 
\textit{clause} is is said to be a disjuction of distinct terms:
$$
t_1 \vee \cdots \vee t_{j_k} = C_k
$$
Then the \textit{satisfiability problem} is the decidability of a conjuction of a set of clauses, 
i.e.:
$$ \wedge_{i=1}^m C_i$$
\end{prob} \cite{skiena2009algorithm}
A \textit{3-SAT problem} is a SAT problem with all clauses having only three boolean variables. 
\begin{definition}[Planar 3-SAT Problem]\label{def:Satisfiability-2}
Given a boolean 3-SAT formula $B$, define the associated graph of $B$ as follows:  
\begin{equation}\label{eqn:sat-1}
G(B) = \left(\set{v_x}{v_x\text{ represents a variable in }B} \cup \set{v_C}{v_C\text{ represent a 
clause in }B}  , \set{\left( v_x, v_C\right) }{x \in C \text{ or } \bar{x} \in C}  \right) 
\end{equation} 
If $G(B)$ in equation (\ref{eqn:sat-1}) is planar, then $B$ is said to be a \textit{Planar 3-SAT 
Problem} \cite{mulzer2008minimum}.
\end{definition} 
\begin{figure}[!h]
\begin{center}
\includegraphics{graphics/LeftSwitchBetweenTwoPolygons.pdf}
\caption{left switch between polygons}
\end{center} 
\end{figure} 
\begin{figure}[!h]
\begin{center}  
\end{center} 
\end{figure} 
\subsection{Logic Engine}
The logic engine simulates the well known Not All Equal 3 SAT Problem (NAE3SAT).  
\subsection{Not All Equal 3 SAT Problem}
\begin{prob}[Satisfiability Problem]\label{prob:Satisfiability-2}%Problem/Question
Give a set of clauses $C$, each containing three boolean variables, can each clause contain at 
least one true variable and one false variable?
\end{prob}
%Universality component
%encoding the engine
\subsubsection{Construction of the Logic Engine}
The components of the logic engine are as follows: the rigid frame, the shaft, the armatures, 
the chains, and the flags.  The \textit{rigid frame} is a rectangular enclosure with a horizontal 
shaft place at mid-height.  The \textit{armatures} are concentric rectangular frames contained 
within the rigid frame.  Each armature can rotate about the shaft; other motions on the armature 
are disallowed.  Given an NAE3SAT, for each variable there is a corresponding armature. On each 
armature, there are chains.  A pair of \textit{chains}, $a_j$ and $\bar{a}_j$ correspond to the 
variable $x_j$ and $\bar{x}_j$ respectively.  The pair is placed on each armature, reflected at a 
height of $h$ above and below the shaft, i.e. one place above the shave at a height of $h$, the 
other placed below the shaft at a height of $-h$.
%insert an armature graphic

\subsection{Encoding the Logic Engine}
For each clause of an NAE3SAT, there exists a set of corresponding chains, namely the $h^\text{th}$ 
clause is the set of chains on the armatures at the $h^\text{th}$ row above and below the shaft. A 
chain is \textit{flagged} if the corresponding variable resides within the clause.  The flag can 
point in either the left or right directions indicating a truth assignment for that variable within 
the clause.  A flag is attached to the $i^\text{th}$ chain of every $a_j^\text{th}$ and 
$\bar{a}_j^\text{th}$ chain with the following exceptions:
\begin{enumerate}
 \item if the variable $x_j$  is in clause $C_i$, then link $i$ of $a_j$ is unflagged,
 \item if the variable $\bar{x}_j$ is in clause $C_i$, then link $i$ of $a_j$ is unflagged.
\end{enumerate}
\begin{thm}\label{thm:Satisfiability-1}
 An instance of $NAE3SAT$ is a ``yes'' instance if and only if the corresponding logic engine has a 
flat, collision-free configuration.
\end{thm}
%add proof
